\documentclass[11pt]{article}
\usepackage[margin=1in]{geometry}
\usepackage{amsmath,amssymb,amsthm}
\usepackage[colorlinks=true,urlcolor=blue,citecolor=blue,linkcolor=blue]{hyperref}
\newtheorem{theorem}{Theorem}
\newtheorem{proposition}[theorem]{Proposition}
\newtheorem{lemma}[theorem]{Lemma}
\newtheorem{corollary}[theorem]{Corollary}
\theoremstyle{remark}
\newtheorem{remark}[theorem]{Remark}
\newcommand{\R}{\mathbb R}
\newcommand{\C}{\mathbb C}
\newcommand{\norm}[1]{\left\lVert#1\right\rVert}
\newcommand{\diag}{\operatorname{diag}}
\newcommand{\rank}{\operatorname{rank}}
\newcommand{\range}{\operatorname{range}}
\newcommand{\MM}{T_{\mathrm{MM}}}
\newcommand{\tO}{\widetilde O}
\title{Product-aware ridge filters:\\a proved primitive and the remaining spectral-cutoff obstacle}
\author{Research working note}
\date{November 14, 2026}
\begin{document}
\maketitle
\begin{abstract}
The desired result remains a gap-free, high-accuracy spectral approximation of
\(AB\) with faster total arithmetic complexity, including its explicit output.
This note does not claim that result or a state-of-the-art improvement for it.
We derive a product-aware ridge-filter primitive from constant-distortion input
sketches, including a uniform operator-error bound, rank-deficient inputs, and
setup costs. We then give exact counterexamples to two tempting extensions and
a coefficient-growth obstruction for uniformly sharp filters assembled from
negative-shift ridge filters. These results isolate a specific missing step;
cheap regularized solves alone do not establish the proposed approximation bound.
The novelty of the primitive itself remains unverified.
\end{abstract}

\section{Objective and cost model}
Let \(A\in\R^{n\times m}\), \(B\in\R^{m\times p}\),
\(m\leq\min(n,p)\), \(M=AB\), \(D=n+p\), and \(I=mD\).
The target in \texttt{research-direction.typ} is to produce width-\(r\) factors
and their explicit product \(C\) such that, for every fixed input,
\[
 \Pr\{\norm{M-C}_2\leq(1+\theta)\sigma_{r+1}(M)\}\geq1-\delta
\]
in construction time \(\tO(Ir+m^{\omega+\zeta})\) and total time including
expansion. Here \(\omega\leq9/4\) is the bound in \cite{amanov}, and
\(\zeta>0\) is fixed exponent slack. The broader absolute-error objective
also permits other algorithms, but any claimed improvement must specify their
input class, tolerance, confidence, rank, and total output cost.

All results below count exact-real arithmetic operations. Matrix inverses in
small cores are algebraic operations; finite-precision stability and bit cost
are not proved here. We use
\(\MM(n,r,p)\leq O(npr)\) for expansion, allowing fast multiplication when
appropriate. A ridge-filter action returns vectors, not an \(n\times n\)
matrix. It does not by itself return the required low-rank approximation.

For the notes' example \(n=p=N\), \(m=\lceil N^{8/9}\rceil\),
\(r=1\), \(\theta=1/m\), the stated exponents check out:
\[
 \omega(1,8/9,1)\leq\frac{625}{288},\qquad
 \omega(8/9,1,8/9)\leq\frac{169}{81},\qquad
 (8/9)(9/4)=2.
\]
These follow from Theorem 1.2 of \cite{amanov}; they are upper bounds for
the named exact-product and Gram routes, not a lower bound against every
possible approximation algorithm. The absolute-error conclusion still needs
\((1+1/m)\sigma_2(M)\leq\varepsilon\).

\section{What ordinary input embeddings already give}
Set
\[
 P=A^\top A,\qquad Q=BB^\top.
\]
Suppose \(S\) and \(T\) satisfy the squared-norm embedding inequalities
\begin{align}
 (1-\eta)P&\preceq (SA)^\top SA\preceq(1+\eta)P,\label{eq:embedding-a}\\
 (1-\eta)Q&\preceq (TB^\top)^\top TB^\top\preceq(1+\eta)Q.
 \label{eq:embedding-b}
\end{align}
The parameter \(\eta\) always denotes distortion of squared norms in this
note. Norm-distortion conventions in references must be converted accordingly.

\begin{proposition}[Lifted sketch approximation]\label{prop:lift}
Let \(U=SA\), \(V=TB^\top\), and let \(\widetilde C_r\) be a truncated
SVD of \(UV^\top\). Then
\[
 C=A U^\dagger\widetilde C_r(V^\dagger)^\top B
 \quad\text{satisfies}\quad
 \rank C\leq r,\qquad
 \norm{M-C}_2\leq\frac{1+\eta}{1-\eta}\sigma_{r+1}(M).
\]
No full-rank assumption is needed.
\end{proposition}
\begin{proof}
The embeddings preserve kernels. Thus \(AU^\dagger U=A\) and
\(V^\top(V^\dagger)^\top B=B\). On the relevant ranges,
\(\norm{AU^\dagger}_2,\norm{B^\top V^\dagger}_2\leq(1-\eta)^{-1/2}\).
The singular vectors retained in \(\widetilde C_r\) lie in \(\range U\)
and \(\range V\). Consequently
\[
 \norm{M-C}_2\leq(1-\eta)^{-1}
       \norm{UV^\top-\widetilde C_r}_2.
\]
The optimal rank-\(r\) truncation \(M_r\) lies in the column and row
spaces of \(M\), hence in those of \(A\) and \(B\). Applying the upper
embedding bounds on these spaces gives
\(\norm{S(M-M_r)T^\top}_2\leq(1+\eta)\norm{M-M_r}_2\).
Optimality of the sketched truncation proves the result.
\end{proof}

For constant \(\eta\), standard fast subspace embeddings have
\(\tO(m)\) rows, can be applied in \(\tO(I)\) time, and leave cores
processable in \(\tO(m^{\omega+\zeta})\) time; see
\cite{cnw,embedding}. A compact core SVD followed by lifting only the
retained coefficient vectors costs \(O(Ir)\), not a full tall basis
construction. This is an existing sketch-and-solve pattern, and is not claimed
as a new algorithm. The issue is the dependence on high accuracy.

\begin{proposition}[Sharpness of distortion-only lifting]\label{prop:reversal}
For \(0<\eta<1\), choose
\(1<\alpha<(1+\eta)/(1-\eta)\). Let
\[
 A=B=\diag(\sqrt\alpha,1),\qquad
 S=T=\diag(\sqrt{1-\eta},\sqrt{1+\eta}).
\]
Both embeddings satisfy \eqref{eq:embedding-a}--\eqref{eq:embedding-b}.
The unique rank-one sketched truncation lifts to
\(C=\diag(0,1)\), with error \(\alpha\), whereas \(\sigma_2(M)=1\).
\end{proposition}
\begin{proof}
The sketched product is
\(\diag((1-\eta)\alpha,1+\eta)\), whose second diagonal entry is larger.
Lifting restores that entry to 1 and does not restore the discarded entry.
\end{proof}

The example pads to arbitrarily tall inputs. It is a counterexample to an
argument using only the two distortion inequalities, not an OSE distributional
lower bound. To obtain the target through Proposition~\ref{prop:lift}, one
needs \(\eta\leq\theta/(2+\theta)\); order-\(\theta\) distortion is also
necessary for a uniform argument of this form. General oblivious embeddings
have the familiar \(\Omega(m/\eta^2)\) row lower bound in sufficiently large
ambient dimension \cite{oselower}. This does not lower-bound adaptive
product-aware algorithms or embeddings using the identity in smaller dimensions.

\section{A uniform ridge-filter primitive}
We now keep \(\eta\) constant. Write
\[
 \widetilde P=(SA)^\top SA,\quad
 \widetilde Q=(TB^\top)^\top TB^\top,\quad
 \mathcal D=\diag(P,Q),\quad
 \widetilde{\mathcal D}=\diag(\widetilde P,\widetilde Q).
\]
For \(\lambda>0\), define the two \(2m\times2m\) block systems
\begin{equation}\label{eq:block}
 \mathcal T_\lambda=
 \begin{pmatrix}P&\sqrt\lambda I_m\\-\sqrt\lambda I_m&Q\end{pmatrix},
 \qquad
 \widetilde{\mathcal T}_\lambda=
 \begin{pmatrix}\widetilde P&\sqrt\lambda I_m\\-\sqrt\lambda I_m&\widetilde Q\end{pmatrix}.
\end{equation}
The coupling is kept exact; only the diagonal Gram blocks are sketched.
Define the ridge filter
\[
 \mathcal R_\lambda=MM^\top(MM^\top+\lambda I_n)^{-1}.
\]

\begin{theorem}[Uniform ridge actions]\label{thm:ridge}
Assume \eqref{eq:embedding-a}--\eqref{eq:embedding-b} with
\(0<\eta<1/2\). For any \(v\in\R^n\), set \(z_0=0\) and perform
\begin{equation}\label{eq:iterate}
 z_{j+1}=z_j+\widetilde{\mathcal T}_\lambda^{-1}
       \left(\binom{A^\top v}{0}-\mathcal T_\lambda z_j\right).
\end{equation}
Let \(x_j\) be the first \(m\) components of \(z_j\), and set
\(\widehat{\mathcal R}_{\lambda,j}v=Ax_j\). Then
\begin{equation}\label{eq:ridge-bound}
 \norm{\widehat{\mathcal R}_{\lambda,j}-\mathcal R_\lambda}_2
 \leq K_\eta\rho_\eta^j,
 \qquad
 K_\eta=\sqrt{\frac{1+\eta}{1-\eta}},\quad
 \rho_\eta=\frac{\eta}{1-\eta}<1.
\end{equation}
The same event establishes this inequality for all positive shifts, all
right-hand sides, and all iteration counts simultaneously, including adaptive
choices. The result includes rank-deficient inputs.
\end{theorem}

\begin{proof}[Proof for positive definite \(P,Q\)]
The symmetric part of \(\mathcal T_\lambda\) is \(\mathcal D\succ0\),
so the system is nonsingular. Likewise the sketched system is nonsingular.
Let
\(\mathcal E=\mathcal D-\widetilde{\mathcal D}
             =\mathcal T_\lambda-\widetilde{\mathcal T}_\lambda\).
The embedding inequalities imply
\[
 \norm{\widetilde{\mathcal D}^{-1/2}\mathcal E
                \widetilde{\mathcal D}^{-1/2}}_2\leq\rho_\eta.
\]
Furthermore,
\[
 \widetilde{\mathcal D}^{-1/2}\widetilde{\mathcal T}_\lambda
       \widetilde{\mathcal D}^{-1/2}=I+J,\qquad J^\top=-J.
\]
For real skew-symmetric \(J\),
\((I+J)^\top(I+J)=I+J^\top J\succeq I\). Hence
\(\norm{(I+J)^{-1}}_2\leq1\), regardless of the size of the coupling.
If \(z_*\) is the exact solution and \(e_j=z_j-z_*\), iteration
\eqref{eq:iterate} gives
\[
 e_{j+1}=-\widetilde{\mathcal T}_\lambda^{-1}\mathcal E e_j,
 \qquad
 \norm{e_{j+1}}_{\widetilde{\mathcal D}}
 \leq\rho_\eta\norm{e_j}_{\widetilde{\mathcal D}}.
\]

Writing \(z_*=(x_*,y_*)\) and \(h=A^\top v\), the second block equation
gives \(y_*=\sqrt\lambda Q^{-1}x_*\), and the first gives
\((P+\lambda Q^{-1})x_*=h\). Thus
\begin{align*}
 \norm{z_*}_{\mathcal D}^2
 &=x_*^\top(P+\lambda Q^{-1})x_*\\
 &=h^\top(P+\lambda Q^{-1})^{-1}h
 \leq h^\top P^{-1}h\leq\norm v_2^2.
\end{align*}
The last inequality uses the orthogonal projector \(AP^{-1}A^\top\).
Since \(z_0=0\),
\[
 \norm{A(x_j-x_*)}_2
 \leq\norm{e_j}_{\mathcal D}
 \leq\frac{\rho_\eta^j}{\sqrt{1-\eta}}
                         \norm{z_*}_{\widetilde{\mathcal D}}
 \leq K_\eta\rho_\eta^j\norm v_2.
\]
Finally, Woodbury's identity gives
\[
 A(P+\lambda Q^{-1})^{-1}A^\top
 =MM^\top(MM^\top+\lambda I_n)^{-1}.
\]
The estimate holds for every \(v\), so it is an operator-norm estimate.
\end{proof}

\begin{proof}[Extension to positive semidefinite \(P,Q\)]
First, both matrices in \eqref{eq:block} are nonsingular even when their
diagonal blocks are singular. For example, if
\(\mathcal T_\lambda(x,y)=0\), then multiplying by \((x^\top,y^\top)\)
gives \(x^\top Px+y^\top Qy=0\). Thus \(Px=Qy=0\), after which
the two block equations give \(x=y=0\).

For the estimate, add \(\tau I_m\) to each of \(P,Q,\widetilde P,
\widetilde Q\), with \(\tau>0\). The same distortion bounds still hold.
Use the same right-hand side \((A^\top v,0)\). The preceding proof gives
\[
 \norm{A(x_{j,\tau}-x_{*,\tau})}_2
 \leq K_\eta\rho_\eta^j\norm v_2,
\]
because \(P+\tau I_m\succeq A^\top A\) and
\(A(P+\tau I_m)^{-1}A^\top\preceq I_n\).
Nonsingularity at \(\tau=0\) implies continuity of every inverse and of
every fixed finite iterate as \(\tau\downarrow0\).
The limiting exact solution satisfies
\[
 x_*=Q(PQ+\lambda I_m)^{-1}A^\top v.
\]
The push-through identity then gives \(Ax_*=\mathcal R_\lambda v\)
without requiring an inverse of \(P\) or \(Q\). Passing to the limit proves
\eqref{eq:ridge-bound}. Rank deficiency therefore adds no regularization
parameter to the actual algorithm.
\end{proof}

\begin{corollary}[Arithmetic accounting]\label{cor:cost}
For constant \(\eta<1/2\), constant-distortion sketches can be constructed
and their Gram matrices formed in \(\tO(I+m^{\omega+\zeta})\) time.
For each distinct requested shift, factoring
\(\widetilde{\mathcal T}_\lambda\) costs
\(O(m^{\omega+\zeta})\), and applying its inverse to \(s\) vectors
costs \(O(m^2s)\) conventionally. Applying the true system uses
\(P x=A^\top(Ax)\) and \(Qy=B(B^\top y)\), costing \(O(Is)\).

Consequently, \(L\) shifts, each applied to a block of \(s\) vectors with
operator error at most \(\xi\in(0,1)\), cost
\[
 \tO\!\left(I+(L+1)m^{\omega+\zeta}
              +IsL\log(1/\xi)\right).
\]
In this expression confidence and dimension logarithms are hidden, while
the displayed refinement logarithm is explicit. There is no polynomial
dependence on a factor condition number or on \(\lambda\).
For polylogarithmically many shifts and \(s=O(r)\), this primitive fits
the desired construction budget up to logarithms. It still lacks a
low-rank approximation theorem. Any eventual \(XY^\top\) must additionally
pay \(\MM(n,r,p)\) to be explicitly formed.
\end{corollary}

The probability statement comes solely from the two input embeddings, each
chosen with failure probability at most \(\delta/2\). Once both hold, the
ridge theorem is deterministic and simultaneous; there is no additional union
bound over adaptively chosen shifts or vectors. The inequality
\(I\geq2m^2\) absorbs dense small-core inverse applications. Inverses can be
precomputed by fast algebra; triangular solves should not silently be charged
\(O(m^2s)\) after a factorization that has not supplied an applicable inverse.

The proof is an application of accretivity and sketch preconditioning, with
connections to classical Hermitian/skew-Hermitian splitting \cite{hss}.
It should not be described as a new AMM algorithm without a further novelty
audit and a proof of the required approximation guarantee.

\section{Why an interior spectral cutoff is different}
Resolvents needed for a sharp spectral cutoff typically involve shifts near
the positive spectrum of \(MM^\top\), rather than negative real shifts.
Keeping the same sketches does not give the preceding contraction there.

\begin{proposition}[Exact failure near a cutoff]\label{prop:shift}
Take scalar Gram blocks \(P=Q=\widetilde Q=1\),
\(\widetilde P=1+\eta\), and replace \(\lambda>0\) in
\eqref{eq:block} by
\(\lambda=-(1+\eta)+i\gamma\), \(\gamma>0\), using any complex
square root. The error iteration in \eqref{eq:iterate} has spectral radius
\(\eta/\gamma\). It diverges when \(\gamma<\eta\).
\end{proposition}
\begin{proof}
Let \(s^2=\lambda\). The sketched system has determinant
\(1+\eta+\lambda=i\gamma\), whereas
\(\mathcal E=\diag(-\eta,0)\). The error matrix
\(-\widetilde{\mathcal T}_\lambda^{-1}\mathcal E\) is rank one and has
nonzero eigenvalue \(\eta/(i\gamma)\).
\end{proof}

This is not a lower bound against other linear-system solvers or different
preconditioners. It invalidates the specific inference that constant Gram
distortion suffices for uniformly cheap interior-shift refinement. In
particular, small-core inversion costs alone cannot justify a rational-filter
runtime if the number of full-input correction steps is unbounded.

\section{Can negative shifts alone build the missing filter?}
Write \(r_\lambda(x)=x/(x+\lambda)\), \(\lambda>0\).
These are the scalar functions supplied by Theorem~\ref{thm:ridge} on
the eigenvalues of \(MM^\top\). The following two elementary bounds
restrict attempts to combine them into a sharp universal filter.

\begin{proposition}[Nonnegative filter combinations]\label{prop:positive}
Suppose
\[
 f(x)=\sum_{j=1}^L a_j\prod_{\ell=1}^{d_j}r_{\lambda_{j\ell}}(x),
 \qquad a_j\geq0,\quad d_j\leq d.
\]
An empty product is 1. For \(b\geq a>0\),
\(f(b)\leq(b/a)^d f(a)\). Thus amplification by \(K>1\) between
1 and \(1+g\) requires
\(d\geq\log K/\log(1+g)\).
\end{proposition}
\begin{proof}
For each factor,
\[
 \frac{r_\lambda(b)}{r_\lambda(a)}
 =\frac ba\frac{a+\lambda}{b+\lambda}\leq\frac ba.
\]
Multiplying and then summing with nonnegative weights proves the assertion.
\end{proof}

\begin{corollary}[A rank-one failure instance for nonnegative filters]
Let \(m\geq65\), \(0<\theta\leq1/8\), and
\(M=\diag(1+2\theta,1,\ldots,1)\). Let \(f\) be nonzero and of the
form in Proposition~\ref{prop:positive}. Draw \(\omega\sim N(0,I_m)\),
set \(v=f(M^\top M)\omega/\norm{f(M^\top M)\omega}_2\), and output
\(C=Mvv^\top\). If
\[
 d\leq\frac{\log((m-1)/16)}{2\log((1+2\theta)^2)},
\]
then \(\norm{M-C}_2>(1+\theta)\sigma_2(M)\) with probability at
least \(5/8\). Thus this strategy needs depth
\(\Omega(\theta^{-1}\log m)\) for uniformly high success probability.
\end{corollary}
\begin{proof}
Put \(s=1+2\theta\) and \(K=f(s^2)/f(1)\). The preceding proposition
and the depth assumption give \(K^2\leq(m-1)/16\).
By Markov's inequality, \(\omega_1^2\leq4\) with probability at least
\(3/4\). The sum of the remaining squares has mean \(m-1\) and variance
\(2(m-1)\), so Chebyshev's inequality bounds the probability that it
is below \((m-1)/2\) by \(8/(m-1)\leq1/8\). The intersection has
probability at least \(5/8\). There, \(v_1^2\leq1/3\), while
\[
 \norm{M(I-vv^\top)}_2^2
 =1+(s^2-1)(1-v_1^2)
 \geq1+\tfrac23(4\theta+4\theta^2)>(1+\theta)^2.
\]
The equality follows by restricting
\(M^\top M=I+(s^2-1)e_1e_1^\top\) to \(v^\perp\).
\end{proof}

This is only a lower bound for the specified algorithm and filter class.
Indeed the signed polynomial \(f(x)=x-1\) recovers the distinguished
direction of this particular instance in one application. The instance
therefore is not a general matvec or product-access lower bound.

Signed coefficients remove that particular argument, but introduce another
quantity that a runtime proof must track.

\begin{theorem}[Coefficient growth for a uniform cutoff]\label{thm:coeff}
Let
\[
 f(z)=\sum_{j=1}^L a_j\prod_{\ell=1}^{d_j}
                         \frac{z}{z+\lambda_{j\ell}},
 \quad a_j\in\C,\quad\lambda_{j\ell}>0,\quad
 \mathcal A=\sum_j|a_j|.
\]
Suppose \(0<g\leq1/8\), \(0<\alpha\leq1/4\),
\[
 \sup_{x\in[1/2,1]}|f(x)|\leq\alpha,
 \qquad |f(1+g)|\geq1.
\]
Then
\begin{equation}\label{eq:coeff}
 \log(4\mathcal A/\alpha)
 \geq\log(4/3)\left(
       \frac{\log(3/(8\alpha))}{\sqrt{8g}}-1\right).
\end{equation}
In particular, along \(g\to0\) and \(\alpha=m^{-c}\) for fixed
\(c>0\), \(\log\mathcal A=\Omega(g^{-1/2}\log m)\).
This holds regardless of the number of factors in each product.
\end{theorem}
\begin{proof}
For \(\operatorname{Re}z\geq0\),
\(|z/(z+\lambda)|\leq1\). Hence \(|f(z)|\leq\mathcal A\) in
that half-plane. In particular this holds on the disk centered at \(3/4\)
of radius \(1/2\). Cauchy's estimate bounds the \(j\)-th Taylor
coefficient in that disk by \(\mathcal A(1/2)^{-j}\).
Every point of \([1/2,1+g]\) is within \(3/8\) of the center.
If \(p_k\) is the degree-\(k\) Taylor truncation, then
\[
 \sup_{x\in[1/2,1+g]}|f(x)-p_k(x)|
 \leq4\mathcal A(3/4)^{k+1}.
\]
Choose \(k=\lceil\log(4\mathcal A/\alpha)/\log(4/3)\rceil\).
The error is at most \(\alpha\), so \(p_k\) is bounded by
\(2\alpha\) on \([1/2,1]\), but
\(|p_k(1+g)|\geq1-\alpha\geq3/4\).

Chebyshev's extremal inequality, after mapping \([1/2,1]\) to
\([-1,1]\) by \(x\mapsto4x-3\), gives
\[
 |p_k(1+g)|\leq2\alpha T_k(1+4g)
 \leq2\alpha\exp(k\sqrt{8g}).
\]
This inequality also holds for complex polynomials: interpolation at the
Chebyshev extrema bounds the absolute value by the sum of absolute Lagrange
weights, which equals \(T_k\) outside the interval. The last step uses
\(T_k(1+u)=\cosh(k\operatorname{arcosh}(1+u))\) and
\(\operatorname{arcosh}(1+u)\leq\sqrt{2u}\).
Therefore \(k\geq\log(3/(8\alpha))/\sqrt{8g}\).
The upper bound \(k\leq\log(4\mathcal A/\alpha)/\log(4/3)+1\)
proves \eqref{eq:coeff}.
\end{proof}

\paragraph{What the coefficient bound does and does not show.}
For a simple linear combination of ridge actions, the usual termwise error
estimate is \(\norm{\widehat f(H)-f(H)}_2\leq\mathcal A\xi\)
when every action is approximated to operator error \(\xi\).
It therefore calls for refinement depth
\(O(\log(\mathcal A/\varepsilon_f))\) to establish final filter
error \(\varepsilon_f\). Products also require tracking their lengths
and intermediate errors. Theorem~\ref{thm:coeff} shows why a claim of
polylogarithmic work cannot simply hide coefficient magnitudes in this
error analysis. It does not prove a computational lower bound: large
coefficients alone do not rule out a different stable evaluation, and
uniform scalar separation on an interval is a sufficient filtering strategy,
not a necessary characterization of every rank-\(r\) approximation algorithm.

\subsection{A sharper degree bound, independent of coefficient size}
The coefficient argument is not needed to obtain the following degree
obstruction. It applies to arbitrary complex coefficients and to a larger pole
region than the negative real axis. Its only approximation-theoretic ingredient
is classical polynomial Chebyshev extremality \cite{extremal}.

\begin{theorem}[Restricted-pole growth]\label{thm:poles}
Let \(a<b<c\) be real. Let \(f=p/q\) be a reduced rational function with
complex coefficients, no pole on \([a,b]\), and \(d=\deg p\).
Suppose every zero \(z\) of \(q\) satisfies
\(\operatorname{Re}z\leq(a+b)/2\). If
\(\sup_{x\in[a,b]}|f(x)|\leq\alpha\), then
\[
 |f(c)|\leq\alpha T_d\!\left(\frac{2c-a-b}{b-a}\right).
\]
The bound is sharp over this class, since constant denominators are allowed.
\end{theorem}
\begin{proof}
For a zero \(z\) of \(q\), convexity of \(|x-z|^2\) and
\(\operatorname{Re}z\leq(a+b)/2\) show that its maximum on \([a,b]\)
is attained at \(b\). Its value is nondecreasing for \(x\geq b\).
Factoring \(q\) over \(\C\) therefore gives
\[
 |q(x)|\leq|q(b)|\quad(a\leq x\leq b),\qquad
 |q(c)|\geq|q(b)|>0.
\]
Thus \(|p(x)|\leq\alpha|q(b)|\) on \([a,b]\). Polynomial Chebyshev
extremality implies
\[
 |p(c)|\leq\alpha|q(b)|
      T_d\!\left(\frac{2c-a-b}{b-a}\right).
\]
For complex coefficients, multiply \(p\) by a unit complex scalar so that
its value at \(c\) is nonnegative real and apply the real-polynomial
inequality to its real part; its degree is at most \(d\). Dividing by
\(|q(c)|\geq|q(b)|\) proves the assertion. Sharpness follows from
\(q=1\) and \(p(x)=\alpha T_d((2x-a-b)/(b-a))\).
\end{proof}

In particular, if \(a=1/2\), \(b=1\), \(c=1+g\),
\(0<\alpha<1\), and \(|f(c)|\geq1\), then
\begin{equation}\label{eq:pole-degree}
 d\geq\frac{\operatorname{arcosh}(1/\alpha)}
                 {\operatorname{arcosh}(1+4g)}
   \geq\frac{\log(1/\alpha)}{\sqrt{8g}}.
\end{equation}
Hence left-half-plane poles, and even poles with real part at most \(3/4\),
give no improvement over the polynomial degree needed for this uniform
top-end separation. A rational filter constructed from \(L\) simple ridge
terms has numerator degree at most \(L\) after a common denominator is
formed, so \eqref{eq:pole-degree} applies regardless of coefficient sizes.

This is a degree bound, not a total-arithmetic lower bound. Many shifted
solves might be processed together by fast rectangular multiplication;
degree alone does not charge a separate full-input pass per pole. Nor does
the theorem cover adaptive filters required to be small only at the actual
tail eigenvalues. The remaining algorithmic issue is whether simultaneous
solves and a stable spectral extraction can exploit batching while retaining
the desired high-accuracy guarantee. No such extraction theorem is proved here.

\section{Literature audit and limits of the present result}
The following records the relevant conclusions from the supplied references
and additional directly related primary sources. It is not an exhaustive
novelty certificate.
\begin{description}
 \item[Halko--Martinsson--Tropp \cite{hmt}.]
 Range finding, power schemes, and a posteriori norm estimates are existing
 tools. A posteriori certification requires its own operator applications;
 it does not turn a loose sketch into a high-accuracy rank-\(r\) solution.
 Relevant parts: Sections 4.3--4.5 and 9--11.
 \item[COD \cite{cod,luo}.]
 The 2017 shrinking argument and the 2021 nuclear/Ky Fan refinement remain
 comparison routes. Their factor-energy dependence must not be replaced by
 a product-only singular-value bound without an additional argument.
 Relevant parts: COD Theorems 2--3; revisited COD Section 3, Theorem 1.
 \item[Block Krylov \cite{musco,chen}.]
 The gap-free spectral guarantee has polynomial inverse-accuracy dependence.
 Chen et al. extend guarantees to intermediate block sizes; this is not a
 result removing that dependence. Relevant parts: Musco--Musco Sections 5--6
 and Appendix A; Chen et al. Sections 3.3--3.5.
 \item[Fast core algebra \cite{demmel}.]
 Fast SVD and related core operations are available reductions, not the
 missing product-interaction sketch. Precision and stability qualifications
 still matter in a bit-cost treatment. Relevant part: Section 6.2.
 \item[Oracle lower bounds \cite{bakshi,zhang}.]
 Matvec lower bounds do not directly lower-bound algorithms given the
 separate factors. The 2026 preprint's theorem concerns a right-projector
 oracle task at constant success probability. It must not be quoted as a
 total-arithmetic lower bound for this input model. Relevant parts:
 Bakshi--Narayanan Theorem 1.1 and Remark 1.3; Zhang et al. Section 3.
 \item[Product-aware sampling \cite{bhojanapalli,wu}.]
 These must be included in a novelty search. Their conditioning and tail
 terms prevent silently substituting the desired gap-free relative spectral
 guarantee. Relevant parts: Bhojanapalli et al. Theorem 3.1 and matrix-product
 application; Wu et al. Theorem 3.1.
 \item[AMM and embeddings \cite{cnw,embedding}.]
 Stable-rank AMM and fast input embeddings are established primitives.
 They support inexpensive constant-distortion cores but do not supply the
 high-accuracy interaction correction used nowhere yet in this note.
 \item[Sketch-powered approximation \cite{sketchpower}.]
 This 2026 preprint already refines low-rank approximation using a larger
 sketch. Its additive spectral-tail term and polynomial accuracy costs need
 to be retained in comparisons. Relevant part: Theorem 2.1 and Algorithm 2.
 \item[Shift-and-invert \cite{garber}.]
 Replacing iteration counts by cheap inverse calls is an established idea.
 Its system-solve costs and gap/accuracy parameters must be accounted for;
 an oracle inverse cannot be treated as free.
\end{description}

\paragraph{Current outcome.}
We have proved a useful uniform product ridge action, not the original AMM
target. Its derivation is self-contained; novelty is unresolved. The lifting
counterexample, interior-shift counterexample, coefficient bound, and
restricted-pole degree bound close specific unjustified shortcuts. None is a
lower bound against the full goal.

\paragraph{Next mathematical issue.}
A successful continuation needs an argument beyond the ridge filter proved
here: for example, a threshold-local preconditioner with
explicitly bounded correction work, or a product-aware rank-selection method
that avoids uniform scalar cutoff separation. Such an argument must include
rank exactly at most \(r\) when claiming the notes' guarantee, failures and
precision, setup and lifting, and explicit expansion. There is currently no
proved candidate in this note that beats the \(N^{169/81+o(1)}\) comparison
in the supplied high-accuracy example.

\begin{thebibliography}{99}
\bibitem{hmt} N. Halko, P.-G. Martinsson, J. Tropp.
\emph{Finding Structure with Randomness} (2011).
\url{https://arxiv.org/abs/0909.4061}.
\bibitem{cod} Y. Mroueh, E. Marcheret, V. Goel.
\emph{Co-Occurring Directions Sketching for Approximate Matrix Multiply} (2017).
\url{https://proceedings.mlr.press/v54/mroueh17a.html}.
\bibitem{luo} L. Luo, C. Chen, G. Xie, H. Ye.
\emph{Revisiting Co-Occurring Directions} (2021).
\url{https://arxiv.org/abs/2009.02553}.
\bibitem{musco} C. Musco, C. Musco.
\emph{Randomized Block Krylov Methods for Stronger and Faster Approximate
Singular Value Decomposition} (2015).
\url{https://arxiv.org/abs/1504.05477}.
\bibitem{demmel} J. Demmel, I. Dumitriu, O. Holtz.
\emph{Fast Linear Algebra Is Stable} (2007).
\url{https://arxiv.org/abs/math/0612264}.
\bibitem{chen} T. Chen, E. Epperly, R. Meyer, C. Musco, A. Rao.
\emph{Does Block Size Matter in Randomized Block Krylov Low-Rank Approximation?}
(2025 preprint).
\url{https://arxiv.org/abs/2508.06486}.
\bibitem{bakshi} A. Bakshi, S. Narayanan.
\emph{Krylov Methods Are (Nearly) Optimal for Low-Rank Approximation} (2023).
\url{https://arxiv.org/abs/2304.03191}.
\bibitem{zhang} H. Zhang et al.
\emph{Matrix-Vector Complexity of Low-Rank Approximation} (2026 preprint).
\url{https://arxiv.org/abs/2609.35840}.
\bibitem{amanov} A. Amanov, M. van den Berg, J. Zuiddam.
\emph{Multiplying Matrices via Polynomials: Exposition and Barriers}
(2026 preprint).
\url{https://arxiv.org/abs/2610.30054}.
\bibitem{cnw} M. Cohen, J. Nelson, D. Woodruff.
\emph{Optimal Approximate Matrix Product in Terms of Stable Rank} (2016).
\url{https://arxiv.org/abs/1507.02268}.
\bibitem{embedding} S. Chenakkod, M. Derezi\'nski, X. Dong.
\emph{Optimal Oblivious Subspace Embeddings with Near-Optimal Sparsity} (2025).
\url{https://arxiv.org/abs/2411.08773}.
\bibitem{oselower} J. Nelson, H. L. Nguyen.
\emph{Lower Bounds for Oblivious Subspace Embeddings} (2014).
\url{https://arxiv.org/abs/1308.3280}.
\bibitem{bhojanapalli} S. Bhojanapalli, P. Jain, S. Sanghavi.
\emph{Tighter Low-Rank Approximation via Sampling the Leveraged Element} (2015).
\url{https://arxiv.org/abs/1410.3886}.
\bibitem{wu} S. Wu, S. Bhojanapalli, S. Sanghavi, A. Dimakis.
\emph{Single Pass PCA of Matrix Products} (2016).
\url{https://arxiv.org/abs/1610.06656}.
\bibitem{sketchpower} S. Chenakkod, M. Derezi\'nski.
\emph{Accelerating Power Method with Fast Sketching for Stronger Low-Rank
Approximation} (2026 preprint).
\url{https://arxiv.org/abs/2605.09755}.
\bibitem{garber} D. Garber et al.
\emph{Faster Eigenvector Computation via Shift-and-Invert Preconditioning} (2016).
\url{https://arxiv.org/abs/1605.08754}.
\bibitem{hss} Z.-Z. Bai, G. Golub, M. Ng.
\emph{Hermitian and Skew-Hermitian Splitting Methods for Non-Hermitian Positive
Definite Linear Systems} (2003).
\url{https://doi.org/10.1137/S0895479801395458}.
\bibitem{extremal} L. Bos, S. Ma'u, S. Waldron.
\emph{Extremal Growth of Polynomials}, Proposition 2.1.
\url{https://www.math.auckland.ac.nz/~waldron/Preprints/Extremal-growth/BosMauWaldron1.pdf}.
\end{thebibliography}
\end{document}
